% =============================================================
% 4.4 การออกแบบองค์ประกอบในเกม
% โครงงาน: การออกแบบและพัฒนาเกมเดินเรือด้วยระบบสั่งการด้วยเสียงเพื่อส่งเสริมทักษะการสื่อสาร
% =============================================================

\section{การออกแบบองค์ประกอบในเกม}

องค์ประกอบภายในเกมได้รับการออกแบบเพื่อสร้างสภาพแวดล้อมจำลองที่สมจริง (Immersive Environment) และเอื้อต่อการเรียนรู้และการฝึกทักษะการสื่อสารภาษาอังกฤษ โดยแบ่งองค์ประกอบออกเป็น 3 ส่วนหลัก ดังนี้

\subsection{การออกแบบตัวละครและบทบาท (Characters and Roles)}

เกมจำลองสถานการณ์การทำงานร่วมกันระหว่างสะพานเดินเรือกับห้องเครื่อง โดยแบ่งบทบาทตัวละครออกเป็น 3 ตำแหน่งหลัก ดังนี้

\begin{enumerate}[label=\arabic*)]
    \item \textbf{กัปตัน:} ตัวแทนของผู้เล่น (Player Character) เป็นผู้มีอำนาจตัดสินใจสูงสุดและเป็นผู้ออกคำสั่งภาษาอังกฤษเพียงผู้เดียว
    \item \textbf{ต้นกลเรือ:} ลูกเรือปัญญาประดิษฐ์ชื่อ Harris ประจำห้องเครื่อง ทำหน้าที่รับคำสั่งเครื่องยนต์และความเร็วจากกัปตันผ่านวิทยุช่องที่ 1 (เช่น ``Slow ahead'' ``Stop engine'' ``Full astern'') แล้วตั้งคันโทรเลขเครื่องยนต์ตามคำสั่ง พร้อมทวนคำสั่ง (Read-back) เพื่อยืนยันความเข้าใจ
    \item \textbf{พลถือท้ายและยามระวังหน้า:} ลูกเรือปัญญาประดิษฐ์ชื่อ Finn ทำหน้าที่รับคำสั่งหางเสือจากกัปตันผ่านวิทยุช่องที่ 2 (เช่น ``Starboard fifteen'' ``Midships'' ``hard-a-port'') เพื่อบังคับหางเสือ พร้อมทวนคำสั่ง และคอยเฝ้าระวังรอบตัวเรือ แจ้งเตือนภัยจากเรดาร์และระยะห่างของสิ่งกีดขวาง เพื่อให้กัปตันนำข้อมูลไปตัดสินใจ
\end{enumerate}

กัปตันเลือกช่องวิทยุเพื่อสั่งการลูกเรือแต่ละคน หากผู้เล่นส่งคำสั่งผิดช่อง เช่น สั่งหางเสือผ่านช่องของต้นกลเรือ ลูกเรือจะแจ้งให้กัปตันสั่งการกับลูกเรือที่รับผิดชอบแทน ซึ่งช่วยฝึกให้ผู้เล่นระบุผู้รับสารให้ถูกต้องตามแบบการสื่อสารบนเรือจริง

\subsection{การออกแบบอุปสรรคและสิ่งกีดขวาง (Obstacles and Hazards)}

เพื่อให้เกมมีความท้าทายและป้องกันไม่ให้ผู้เล่นท่องจำเส้นทาง โหมดเดินเรืออิสระจึงสุ่มวางสิ่งกีดขวางใหม่ทุกครั้ง ส่วนภารกิจในโหมดเนื้อเรื่องใช้แผนที่คงที่ ร่องน้ำมีชายฝั่งและแนวหินกั้นทั้งสองข้าง และมีสิ่งกีดขวางกระจายอยู่ภายในร่องน้ำ ดังนี้

\begin{itemize}
    \item \textbf{หินและแนวหินโสโครก:} กำหนดให้ผู้เล่นต้องประเมินระยะห่างและสั่งการเลี้ยวหลบให้ทันเวลา การชนแต่ละครั้งทำให้ความทนทานของตัวเรือลดลงและเสียตู้สินค้า 1 ตู้
    \item \textbf{ทุ่นระเบิด:} ปรากฏตั้งแต่งานระดับ Normal ขึ้นไป เป็นกลุ่มละ 2--3 ลูก การชนทุ่นระเบิดทำให้เสียตู้สินค้า 2 ตู้
    \item \textbf{สันดอน:} ทำให้เรือเกยตื้นโดยตัวเรือไม่เสียหาย แต่ผู้เล่นต้องสั่งถอยหรือเดินหน้าช้าเพื่อนำเรือออก
\end{itemize}

\subsection{การออกแบบกราฟิกและเสียงประกอบ (Graphic and Audio Design)}

การออกแบบด้านทัศนศิลป์และเสียงมุ่งสร้างความรู้สึกดื่มด่ำ ให้ผู้เล่นรู้สึกราวกับบัญชาการเรือจริง

\begin{itemize}
    \item \textbf{กราฟิก (Graphic):} นำเสนอในมุมมองไอโซเมตริก เพื่อให้ผู้เล่นมองเห็นตัวเรือ สภาพแวดล้อมทางทะเล และตำแหน่งของสิ่งกีดขวางโดยรอบได้อย่างชัดเจนในมุมกว้าง
    \item \textbf{เสียงแวดล้อม (Ambient Audio):} เสียงคลื่นและนกนางนวลสังเคราะห์ด้วยโค้ดขณะเล่นจึงไม่วนซ้ำ ร่วมกับเสียงเครื่องยนต์เรือที่ดังต่อเนื่อง เพื่อสร้างบรรยากาศการเดินเรือที่สมจริง
    \item \textbf{เสียงเอฟเฟกต์:} ใช้เสียงสัญญาณเตือนภัย (Alarm) จากเรดาร์เมื่อเรือเคลื่อนเข้าใกล้สิ่งกีดขวาง เพื่อกระตุ้นความตื่นตัวและสร้างความกดดันด้านเวลาให้ผู้เล่นต้องรีบออกคำสั่ง
\end{itemize}